\section{Related Work}
\label{sec:related}

\paragraph{Document parsing pipelines.}
Classical document AI decomposes the task into layout analysis, OCR, and per-region specialist recognition. Modern open-source exemplars are MinerU~\cite{mineru}, which chains DocLayout-YOLO, a formula detector, UniMERNet and a table recognizer; Marker~\cite{marker}, built around the Surya OCR family (a $\sim$650M-parameter recognition VLM); Docling~\cite{docling}; and PP-Structure\,V3~\cite{ppstructurev3} from the PaddleOCR ecosystem. These systems are CPU-\emph{capable} in principle but are engineered and evaluated with GPUs; on CPU we measure PP-StructureV3's server profile at 58\,s/page and 8\,GB RAM (\cref{sec:experiments}). PRISM belongs to this pipeline family but is, to our knowledge, the first designed \emph{from the start} for a strict CPU/memory/size budget and evaluated end-to-end on the full official benchmark under it.

\paragraph{Vision--language models for documents.}
Specialized document VLMs now top every benchmark: PaddleOCR-VL (0.9B)~\cite{paddleocrvl}, MinerU2.5 (1.2B)~\cite{mineru25}, GLM-OCR, dots.ocr, DeepSeek-OCR~\cite{deepseekocr}, and general VLMs such as Gemini and Qwen-VL. Their accuracy is bought with billions of multiply--accumulates per output token and GPU inference. The ``efficient'' end of this family --- SmolDocling-256M~\cite{smoldocling} and GraniteDocling-258M --- matches PRISM's \emph{size} class but not its efficiency class: as autoregressive generators they still emit the entire page token by token, which on CPU costs 55--92\,s per page (\cref{tab:cpu_frontier}), and their accuracy on OmniDocBench-style evaluation trails classical pipelines. Nougat~\cite{nougat} and GOT-OCR2.0~\cite{got} are earlier single-model parsers with similar deployment profiles.

\paragraph{Component models.}
PRISM builds on strong small components from the open ecosystem: PP-DocLayout~\cite{ppdoclayout} (RT-DETR~\cite{rtdetr} layout detection), PP-OCR~\cite{ppocr} text detection/recognition via RapidOCR ONNX runtimes, the Texo distillation~\cite{texo} of formula recognition into 20M parameters, Table Transformer (TATR)~\cite{tatr}, and SLANet-plus~\cite{slanet} for table structure. Our contribution is not these models but the demonstration that the \emph{orchestration layer} --- routing, geometric post-processing, and evaluation-aware output normalization --- closes most of the gap to systems tens to thousands of times larger.

\paragraph{Benchmarks.}
OmniDocBench~\cite{omnidocbench} is the current standard for end-to-end document parsing: nine document genres, bilingual, with component-wise metrics (normalized edit distance for text and reading order, CDM~\cite{cdm} for formulas, TEDS~\cite{teds} for tables) and an official matching-based harness. We use v1.6 (1651 pages) as our primary target and additionally report the v1.5 cut for comparability with published rows. Fox~\cite{fox} and DocGenome are complementary; we focus on OmniDocBench as it is the only benchmark with a maintained public leaderboard across both VLMs and pipelines.
